
\section{Izobarik jarayon}

\begin{enumerate}

\item Izobarik jarayon deb qanday jarayonga aytiladi?
\begin{enumerate}[label=\Alph*)]
\item doimiy bosim ostida kechuvchi jarayonga
\item doimiy haroratda kechuvchi jarayonga
\item issiqlik almashmasdan kechuvchi jarayonga
\item doimiy hajmda kechuvchi jarayonga
\end{enumerate}

\item Izobarik jarayonda ...
\begin{enumerate}[label=\Alph*)]
\item $p$ va $V$ o‘zgaradi, $T$ o‘zgarmaydi.
\item $p$ va $T$ o‘zgaradi, $V$ o‘zgarmaydi.
\item $V$ va $T$ o‘zgaradi, $p$ o‘zgarmaydi.
\item $p, V$ va $T$ o‘zgaradi va tashqi muhit bilan issiqlik almashinish bo‘lmaydi.
\end{enumerate}

\item Agar ideal gaz 270 K ga izobarik ravishda sovitilganda, hajmi 4 marta kamaygan bo‘lsa, uning dastlabki temperaturasi necha $^\circ C$ bo‘lgan?
\begin{enumerate}[label=\Alph*)]
\item 60
\item 81
\item 87
\item 807
\end{enumerate}

\item O‘txonadan trubaga kirgan gazlar $1320^\circ C$ dan $200^\circ C$ temperaturagacha soviydi. Ularning hajmi necha marta kamayadi?
\begin{enumerate}[label=\Alph*)]
\item 1,5
\item 3,4
\item 2,1
\item 4
\end{enumerate}

\item Agar gaz $27^\circ C$ da 6 l hajmga ega bo‘lsa, $77^\circ C$ da qanday hajmni egallaydi (l)?
\begin{enumerate}[label=\Alph*)]
\item 2
\item 14
\item 5
\item 7
\end{enumerate}

\item Gaz 12,32 l hajmni egallagan. U o‘zgarmas bosimda 45 K ga sovitildi, uning hajmi 10,52 l bo‘lib qoldi. Gazning boshlang‘ich temperaturasi qanday bo‘lgan (K)?
\begin{enumerate}[label=\Alph*)]
\item 300
\item 265
\item 308
\item 900
\end{enumerate}

\item Gaz $27^\circ C$ dan $39^\circ C$ gacha isitildi. Agar bosim o‘zgarmagan bo‘lsa, gaz hajmi necha \% ortgan?
\begin{enumerate}[label=\Alph*)]
\item 4
\item 30,7
\item 44,4
\item 50
\end{enumerate}

\item O‘zgarmas bosimda gazni $0^\circ C$ dan necha $^\circ C$ gacha isitganda uning hajmi 2 marta oshadi?
\begin{enumerate}[label=\Alph*)]
\item 100
\item 173
\item 120
\item 273
\end{enumerate}

\item Qanday temperaturali (C) gaz $1^\circ C$ isitilganda hajmi 2 marta ortadi? ($p=const$)
\begin{enumerate}[label=\Alph*)]
\item 0
\item 1
\item -272,15
\item 272
\end{enumerate}

\item Gaz $27^\circ C$ dan $57^\circ C$ gacha isitilgan. Agar bosim o‘zgarmagan bo‘lsa, gaz hajmi necha protsent ortgan?
\begin{enumerate}[label=\Alph*)]
\item 4
\item 8
\item 10
\item 20
\end{enumerate}

\item Gaz o‘zgarmas bosimda $27^\circ C$ dan $51^\circ C$ gacha isitilganda hajmi necha foizga ortadi?
\begin{enumerate}[label=\Alph*)]
\item 7
\item 9
\item 10
\item 8
\end{enumerate}

\item Gaz temperaturasi 60 K ga ortganda uning hajmi 1 l ga ortdi. Agar temperatura yana 30 K ga ortsa, hajmi dastlabki hajmga qaraganda qanchaga ortadi (l)? ($p=const$)
\begin{enumerate}[label=\Alph*)]
\item 2,5
\item 1,5
\item 3,8
\item 0,15
\end{enumerate}

\item Gazni o‘zgarmas bosimda $1^\circ C$ isitganda uning hajmi ikki marta ortgan. Isitish qanday temperaturalar intervalida sodir bo‘lgan (K)?
\begin{enumerate}[label=\Alph*)]
\item (4;6)
\item (2;4)
\item (-2;0)
\item (0,5;2)
\end{enumerate}

\item Ideal gaz hajmining termik koeffitsiyenti nimaga teng?
\begin{enumerate}[label=\Alph*)]
\item $\Delta p / p_0$
\item $-p / (p_0 T)$
\item $\Delta V / (V_0 \Delta T)$
\item $\Delta V / V_0$
\end{enumerate}

\item Hajmi $250 \, sm^3$ bo‘lgan og‘zi ochiq shisha kolba $127^\circ C$ gacha qizdirildi va ochiq og‘zi bilan suvga tushirildi. Natijada kolba $7^\circ C$ temperaturagacha sovidi. Bunda kolbaga necha gramm suv kirgan?
\begin{enumerate}[label=\Alph*)]
\item 125
\item 75
\item 50
\item 25
\end{enumerate}

\item Agar ochiq idishdagi ideal gazning absolyut temperaturasi 1,1 marta ortsa, gaz molekulalarining konsentratsiyasi qanday o‘zgaradi?
\begin{enumerate}[label=\Alph*)]
\item 1,1 marta ortadi
\item 1,1 marta kamayadi
\item 1,21 marta kamayadi
\item o‘zgarmaydi
\end{enumerate}

\item Hajmi $0,1 \, m^3$ bo‘lgan silindrik idish yuzi $0,01 \, m^2$ bo‘lgan porshen bilan yopilgan. Gazning absolyut temperaturasi izobarik ravishda 4 marta pasaytirilganda porshen idish tubidan necha metr balandlikda bo‘ladi?
\begin{enumerate}[label=\Alph*)]
\item 1,52
\item 2
\item 2,5
\item 3
\end{enumerate}

\end{enumerate}



\begin{table}[h!]
\centering
\begin{tabular}{|c|c|c|c|c|c|c|c|c|c|}
\hline
№ & J. & № & J. & № & J. & № & J. & № & J. \\ \hline
1 & A & 5 & D & 9 & C & 13 & D & 17 & C \\ \hline
2 & C & 6 & C & 10 & C & 14 & D & & \\ \hline
3 & C & 7 & A & 11 & D & 15 & B & & \\ \hline
4 & B & 8 & D & 12 & B & 16 & B & & \\ \hline
\end{tabular}
\end{table}